\section{Related work and interpretation}
\label{app:related}

Predictive feature learning now spans several distinct settings. I-JEPA and
V-JEPA learn from masked image and video content
\citep{assran2023ijepa,bardes2024vjepa}; Image World Models study prediction
of photometric transformations and the roles of conditioning and predictor
capacity \citep{garrido2024iwm}. V-JEPA~2 connects video pretraining to
action-conditioned robot planning \citep{assran2025vjepa2}, while
V-JEPA~2.1 adds dense and intermediate-layer supervision to improve spatial
and temporal features \citep{murlabadia2026vjepa21}. These successes motivate
testing which task distinctions remain accessible after prediction in
feature space, especially when they occupy little of the image.

Latent planning can use either pretrained visual features or representations
learned jointly with dynamics. DINO-WM predicts spatial patch features from
a pretrained DINOv2 encoder \citep{zhou2025dinowm,oquab2023dinov2}.
PLDM learns predictive dynamics from reward-free offline interactions
\citep{sobal2025latentplanning}; LeWorldModel and LpWM study end-to-end
learning with different distribution regularizers
\citep{maes2026leworldmodel,kuang2026lpwm}. Our positive control belongs to
the latter line, with its particular recipe stated in
Appendix~\ref{app:planning}. Its stronger planning motivates a controlled
comparison of learned cue binding; it does not establish that sparsity alone
explains the difference.

Regularization and identifiability provide complementary perspectives.
VICReg explicitly controls feature variance and covariance
\citep{bardes2022vicreg}; LeJEPA studies isotropic Gaussian regularization
\citep{balestriero2025lejepa}, and Rectified LpJEPA develops sparse,
nonnegative alternatives \citep{kuang2026rectified}.
In \emph{When Does LeJEPA Learn a World Model?}, Klindt et al.
\citep{klindt2026when} establish linear recovery under assumptions including
independent Gaussian latent variables and stationary additive-noise
transitions. This is an encoder-side result under a specified data model,
not a general guarantee for an action-conditioned predictor. We do not
establish those assumptions for Tunnel Rooms, so our findings do not
contradict that theorem.

Feature selection can nevertheless favor the wrong information for a task.
Sobal et al. \citep{sobal2022slow} show that JEPAs can prefer stationary
noise over a moving object in controlled experiments. Littwin et al.
\citep{littwin2024jepa} analyze a bias toward more predictable features in
deep linear self-distillation networks. These results motivate examining
the relation between predictability and usefulness, but neither identifies
the cause of our cue-binding gap: our cues persist and determine destination
pairings. The transition-exposure experiment therefore tests a separate
possibility and remains inconclusive.

Additional training signals offer ways to align representations with action.
Sensorimotor World Models use inverse dynamics
\citep{ivashkov2026sensorimotor}; auxiliary-task theory explains how added
targets can preserve otherwise unresolved distinctions \citep{yu2025auxjepa}.
Value-guided JEPA planning explicitly aligns latent distance with a
goal-reaching value function \citep{destrade2025value}, while PSG-JEPA adds
training-time physical-state grounding \citep{yan2026psgjepa}.
Temporal Straightening regularizes latent trajectories for planning
\citep{wang2026temporal}; its learned feature aggregation inspires our
different spatial readout (Appendix~\ref{app:aggregator}). Together, these
directions motivate testing which additional constraints help retain cue
identity and destination, beyond improving average prediction or marker
visibility.
